% GENERATED FILE. Edit canonical lessons, then run npm run generate:books.
% canonical-source-hash: fnv1a64:5746977f6cd26140
% canonical-lessons: GE-C302-rente, GE-C302-ruhestand, GE-C302-rentner, GE-C302-jubilaeum, GE-C302-abschied

\chapter{Der Ruhestand}
\label{ch:ge-der-ruhestand}
\hlchaptermodality{302}

\begin{chapteropening}
\textbf{By the end of this chapter:} \emph{I can tell about later life: a pension, retirement, a pensioner, an anniversary and a farewell.}

Complete the last lesson of chapter 302: I can tell about later life: a pension, retirement, a pensioner, an anniversary and a farewell.
\end{chapteropening}

\section[die Rente]{die Rente — the pension}

\label{lesson:GE-C302-rente}



\begin{quote}
Before the new one: say the German for the promotion, then the German for the turning point.
\end{quote}

\subsection*{You'll want to know: die Rente}
\textbf{die Rente} — \textquotedblleft{}the pension\textquotedblright{}.

\begin{quote}
Mit fünfundsechzig ging mein Vater in Rente.
\end{quote}

\begin{quote}
\emph{At sixty-five my father retired.}
\end{quote}

\emph{In Rente gehen} — to retire. In Austria: \emph{die Pension}.

\begin{cousinweb}
\emph{Rente} was borrowed from French \emph{rente} \textquotedblleft{}income\textquotedblright{}, which is usually traced to Latin \emph{reddere} \textquotedblleft{}to give back\textquotedblright{} (Lewis \& Short); EtymDB lists no parent.
\end{cousinweb}

\begin{grammarlens}[title={What you've built}]
One more word for telling, explaining and describing at B1.
\end{grammarlens}

\subsection*{Your turn}
\begin{itemize}
\raggedright
  \item \emph{Say it:} \emph{die Rente}
  \item \emph{Say it:} \emph{die Rente}, once more
  \item \emph{Say it:} a sentence of your own with \emph{die Rente}
  \item \emph{Recall:} say the German for the promotion, then the German for the turning point, then say \emph{die Rente} again
\end{itemize}

\begin{tcolorbox}[breakable,colback=teal!4,colframe=teal!35!black,title={Before you move on}]
What does die Rente mean? (\textquotedblleft{}The pension\textquotedblright{}.) Say it once more.
\end{tcolorbox}
\section[der Ruhestand]{der Ruhestand — the retirement}

\label{lesson:GE-C302-ruhestand}



\begin{quote}
Before the new one: say the German for the turning point, then the German for the pension.
\end{quote}

\subsection*{You'll want to know: der Ruhestand}
\textbf{der Ruhestand} — \textquotedblleft{}the retirement\textquotedblright{}.

\begin{quote}
Im Ruhestand reiste sie viel.
\end{quote}

\begin{quote}
\emph{In retirement she travelled a lot.}
\end{quote}

Formal. \emph{In den Ruhestand gehen} — to retire.

\begin{cousinweb}
\emph{Ruhestand} joins \emph{Ruhe} \textquotedblleft{}rest\textquotedblright{}, from Proto-Germanic \emph{rōwō} \textquotedblleft{}quiet, rest\textquotedblright{} (EtymDB), and \emph{Stand} \textquotedblleft{}standing, state\textquotedblright{}: the state of rest.
\end{cousinweb}

\begin{grammarlens}[title={What you've built}]
One more word for telling, explaining and describing at B1.
\end{grammarlens}

\subsection*{Your turn}
\begin{itemize}
\raggedright
  \item \emph{Say it:} \emph{der Ruhestand}
  \item \emph{Say it:} \emph{der Ruhestand}, once more
  \item \emph{Say it:} a sentence of your own with \emph{der Ruhestand}
  \item \emph{Recall:} say the German for the turning point, then the German for the pension, then say \emph{der Ruhestand} again
\end{itemize}

\begin{tcolorbox}[breakable,colback=teal!4,colframe=teal!35!black,title={Before you move on}]
What does der Ruhestand mean? (\textquotedblleft{}The retirement\textquotedblright{}.) Say it once more.
\end{tcolorbox}
\section[der Rentner]{der Rentner — the pensioner}

\label{lesson:GE-C302-rentner}



\begin{quote}
Before the new one: say the German for the pension, then the German for the retirement.
\end{quote}

\subsection*{You'll want to know: der Rentner}
\textbf{der Rentner} — \textquotedblleft{}the pensioner\textquotedblright{}.

\begin{quote}
Viele Rentner sitzen am Morgen im Park.
\end{quote}

\begin{quote}
\emph{Many pensioners sit in the park in the morning.}
\end{quote}

Plural: \emph{die Rentner}. Feminine \emph{die Rentnerin}.

\begin{cousinweb}
\emph{Rentner} is built on \emph{Rente} \textquotedblleft{}pension\textquotedblright{}, borrowed from French \emph{rente} and usually traced to Latin \emph{reddere} \textquotedblleft{}to give back\textquotedblright{} (Lewis \& Short).
\end{cousinweb}

\begin{grammarlens}[title={What you've built}]
One more word for telling, explaining and describing at B1.
\end{grammarlens}

\subsection*{Your turn}
\begin{itemize}
\raggedright
  \item \emph{Say it:} \emph{der Rentner}
  \item \emph{Say it:} \emph{der Rentner}, once more
  \item \emph{Say it:} a sentence of your own with \emph{der Rentner}
  \item \emph{Recall:} say the German for the pension, then the German for the retirement, then say \emph{der Rentner} again
\end{itemize}

\begin{tcolorbox}[breakable,colback=teal!4,colframe=teal!35!black,title={Before you move on}]
What does der Rentner mean? (\textquotedblleft{}The pensioner\textquotedblright{}.) Say it once more.
\end{tcolorbox}
\section[das Jubiläum]{das Jubiläum — the anniversary, the jubilee}

\label{lesson:GE-C302-jubilaeum}



\begin{quote}
Before the new one: say the German for the retirement, then the German for the pensioner.
\end{quote}

\subsection*{You'll want to know: das Jubiläum}
\textbf{das Jubiläum} — \textquotedblleft{}the anniversary, the jubilee\textquotedblright{}.

\begin{quote}
Die Firma feierte ihr hundertjähriges Jubiläum.
\end{quote}

\begin{quote}
\emph{The company celebrated its hundredth anniversary.}
\end{quote}

Plural: \emph{die Jubiläen}. For a firm, a club or years at work.

\begin{cousinweb}
\emph{Jubiläum} is a learned borrowing from Late Latin \emph{iūbilaeus} \textquotedblleft{}of the jubilee\textquotedblright{}, which came from Hebrew through Greek; EtymDB lists no parent, so this follows the usual account.
\end{cousinweb}

\begin{grammarlens}[title={What you've built}]
One more word for telling, explaining and describing at B1.
\end{grammarlens}

\subsection*{Your turn}
\begin{itemize}
\raggedright
  \item \emph{Say it:} \emph{das Jubiläum}
  \item \emph{Say it:} \emph{das Jubiläum}, once more
  \item \emph{Say it:} a sentence of your own with \emph{das Jubiläum}
  \item \emph{Recall:} say the German for the retirement, then the German for the pensioner, then say \emph{das Jubiläum} again
\end{itemize}

\begin{tcolorbox}[breakable,colback=teal!4,colframe=teal!35!black,title={Before you move on}]
What does das Jubiläum mean? (\textquotedblleft{}The anniversary, the jubilee\textquotedblright{}.) Say it once more.
\end{tcolorbox}
\section[der Abschied]{der Abschied — the farewell, the parting}

\label{lesson:GE-C302-abschied}



\begin{quote}
Before the new one: say the German for the pensioner, then the German for the anniversary.
\end{quote}

\subsection*{You'll want to know: der Abschied}
\textbf{der Abschied} — \textquotedblleft{}the farewell, the parting\textquotedblright{}.

\begin{quote}
Beim Abschied weinten alle.
\end{quote}

\begin{quote}
\emph{At the farewell everyone cried.}
\end{quote}

Plural: \emph{die Abschiede}. \emph{Abschied nehmen} — to say goodbye.

\begin{cousinweb}
\emph{Abschied} is built on \emph{abscheiden} \textquotedblleft{}to depart\textquotedblright{}, from \emph{scheiden} \textquotedblleft{}to separate\textquotedblright{}, Proto-Germanic \emph{skaiþaną} (EtymDB).
\end{cousinweb}

\begin{grammarlens}[title={What you've built}]
One more word for telling, explaining and describing at B1.
\end{grammarlens}

\subsection*{Your turn}
\begin{itemize}
\raggedright
  \item \emph{Say it:} \emph{der Abschied}
  \item \emph{Say it:} \emph{der Abschied}, once more
  \item \emph{Say it:} a sentence of your own with \emph{der Abschied}
  \item \emph{Recall:} say the German for the pensioner, then the German for the anniversary, then say \emph{der Abschied} again
\end{itemize}

\begin{tcolorbox}[breakable,colback=teal!4,colframe=teal!35!black,title={Before you move on}]
What does der Abschied mean? (\textquotedblleft{}The farewell, the parting\textquotedblright{}.) Say it once more.
\end{tcolorbox}
